\chapter{SABR in Rates and the Volatility Cube}\label{ch:rc:sabr-in-rates-and-the-volatility-cube}

In 2015 a euro rates desk asked its smile model for the volatility of a
receiver swaption struck at minus twenty basis points. The model was SABR, the
market's standard since 2002, in the lognormal form every desk used: its
formula takes $\ln(F/K)$ and $(FK)^{(1-\beta)/2}$, and with a negative strike it
had nothing to return. Within months the whole euro market had moved to one of
two fixes: add a constant to every rate before applying the formula, or use the
version of the model in which rates are normal. Each fix fits the quoted
strikes; they disagree about the strikes nobody quotes. This chapter applies
SABR (One Quant Book 5, chapter 11) to rates, builds the swaption volatility
cube from it smile by smile, interpolates it, and looks where it breaks: at
negative strikes, and in the low-strike wing of long expiries, where the
formula's density turns negative.

\section{SABR in rates: normal SABR and the backbone}

A forward swap rate $F$ under its annuity measure has SABR dynamics
\[
dF_t = \alpha_t F_t^{\beta}\,dW_t,\qquad d\alpha_t = \nu\alpha_t\,dZ_t,\qquad
d\langle W,Z\rangle_t = \rho\,dt ,
\]
with the parameters $(\alpha,\beta,\rho,\nu)$ of Book 5. Rates desks fix
$\beta$ by convention rather than by fit, since $\beta$ and $\rho$ both tilt the
smile and a single section cannot separate them; the choice sets the
backbone, how the at-the-money volatility moves when the forward moves.

\begin{definition}[Normal SABR]\label{def:rc:sabr-in-rates-and-the-volatility-cube:normal}
\emph{Normal SABR}\index{normal SABR} is SABR with $\beta=0$: $dF_t=\alpha_t\,dW_t$,
a rate that is normal given its volatility and may take any sign. Its
at-the-money normal volatility does not depend on the level of the forward:
the backbone is flat in normal terms.
\end{definition}

\begin{proposition}[The normal-volatility expansion of normal SABR]\label{prop:rc:sabr-in-rates-and-the-volatility-cube:normalexp}
To first order in the expiry $T$ (Hagan et al., 2002), the normal implied
volatility of \omterm{def:rc:sabr-in-rates-and-the-volatility-cube:normal}{normal SABR} is
\[
\sigma_N(K) = \alpha\,\frac{\zeta}{x(\zeta)}\Bigl(1+\frac{2-3\rho^2}{24}\,\nu^2T\Bigr),
\quad \zeta = \frac{\nu}{\alpha}(F-K),\quad
x(\zeta) = \ln\frac{\sqrt{1-2\rho\zeta+\zeta^2}+\zeta-\rho}{1-\rho},
\]
with $\zeta/x(\zeta)=1$ at the money. The smile is a function of $F-K$ only:
it moves with the forward, strike for strike in basis points.
\end{proposition}
\admitted\ (Hagan et al., 2002, derive it by singular perturbation; One Quant
Book 5, chapter 11, gives the lognormal case.)

\omcode{code/firm/sabrcube/firm_sabrcube.py}{41}{61}{Hagan's lognormal expansion on shifted rates, and the normal-SABR normal volatility.}\label{lst:rc:sabr-in-rates-and-the-volatility-cube:hagan}

\section{Negative rates and shifted models}

\begin{definition}[Shifted lognormal model, shifted SABR]\label{def:rc:sabr-in-rates-and-the-volatility-cube:shifted}
A \emph{shifted lognormal model}\index{shifted lognormal model} with shift
$\zeta>0$ makes $F+\zeta$ lognormal, $dF_t=\sigma(F_t+\zeta)\,dW_t$, so that the
rate is bounded below by $-\zeta$; options are priced with Black's formula on
$F+\zeta$ and $K+\zeta$. \emph{Shifted SABR}\index{shifted SABR} applies SABR's
dynamics to $F+\zeta$, $dF_t=\alpha_t(F_t+\zeta)^\beta\,dW_t$, and Hagan's lognormal
expansion to the shifted forward and strike.
\end{definition}

The shift is a floor on the modelled rate, chosen in advance. It should sit
below any rate the market can plausibly reach.

\begin{definition}[Effective lower bound]\label{def:rc:sabr-in-rates-and-the-volatility-cube:elb}
The \emph{effective lower bound}\index{effective lower bound} of a policy rate is
the level below which a central bank is not expected to cut, because further
cuts would push cash out of the banking system or hurt banks more than they
ease conditions. Rates models use it to bound the rate distribution, and a
shifted model's shift must exceed it in magnitude.
\end{definition}

\begin{dated}{2026-09}{Where the lower bound has been}\label{dat:rc:sabr-in-rates-and-the-volatility-cube:elb}
The ECB's deposit rate was negative from June 2014 to July 2022, with a low of
$-0.50\%$ from September 2019; the Swiss National Bank's policy rate was
$-0.75\%$ from January 2015 and returned above zero in September 2022; the Bank
of Japan's was $-0.1\%$ from 2016 until 19 March 2024.
\end{dated}

\begin{example}[One smile, three models]\label{ex:rc:sabr-in-rates-and-the-volatility-cube:neg}
A one-year option on a two-year euro swap, forward $-0.30\%$, is quoted at
normal volatilities of 29.3, 26.5, 25.7, 27.8, 31.4 and 39.8 basis points at
strikes 50 and 25 basis points below the forward, at the money, and 25, 50 and
100 above (a smile generated by \omterm{def:rc:sabr-in-rates-and-the-volatility-cube:normal}{normal SABR}, illustrative). \omterm{def:rc:sabr-in-rates-and-the-volatility-cube:shifted}{Shifted SABR} with
$\beta=0.5$ fits it to 0.77 basis points of root-mean-square error with a shift
of 1\%, to 0.11 with a shift of 3\%; \omterm{def:rc:sabr-in-rates-and-the-volatility-cube:normal}{normal SABR} fits it exactly. Asked for a
receiver struck at $-0.90\%$, sixty basis points below the forward and outside
the quotes, they answer 0.177, 0.280 and 0.295 basis points of annuity
(\cref{fig:rc:sabr-in-rates-and-the-volatility-cube:neg}): with a 1\% shift the
modelled rate cannot fall below $-1\%$, and the deep receiver is 37\% cheaper
than with 3\%.
\end{example}

\begin{omfigure}
\begin{tikzpicture}
\begin{axis}[width=12cm,height=5.8cm,
  xlabel={strike (\%)},ylabel={normal volatility (bp)},
  tick label style={font=\small},label style={font=\small},
  xmin=-1,xmax=0.8,ymin=24,ymax=45,grid=major,grid style={gray!25},unbounded coords=jump,
  legend columns=4,legend style={font=\footnotesize,at={(0.5,-0.25)},anchor=north,draw=none,fill=none,/tikz/every even column/.append style={column sep=6pt}},
  table/col sep=comma]
\addplot[omThm,thick] table[x=k,y=s1]{figdata/rates-credit-risk/05-sabr-in-rates-and-the-volatility-cube/negsmile.csv};
\addplot[omDef,thick,dashed] table[x=k,y=s3]{figdata/rates-credit-risk/05-sabr-in-rates-and-the-volatility-cube/negsmile.csv};
\addplot[omProp,thick,dotted] table[x=k,y=nrm]{figdata/rates-credit-risk/05-sabr-in-rates-and-the-volatility-cube/negsmile.csv};
\addplot[only marks,mark=*,mark size=2pt,black] table[x=k,y=vol]{figdata/rates-credit-risk/05-sabr-in-rates-and-the-volatility-cube/negquotes.csv};
\legend{shifted SABR 1\%,shifted SABR 3\%,normal SABR,quotes}
\draw[gray,dashed] (axis cs:-0.9,24) -- (axis cs:-0.9,45) node[pos=0.93,right,font=\footnotesize,text=black] {$-0.90\%$};
\end{axis}
\end{tikzpicture}

\omcaption{A negative-rate smile (one year into two years, forward $-0.30\%$)
fitted by three models, in normal volatility. The 3\% shift and \omterm{def:rc:sabr-in-rates-and-the-volatility-cube:normal}{normal SABR}
pass through the quotes; the 1\% shift misses them by up to 1.2 basis points and, below them, its
volatility falls towards its floor at $-1\%$. Data: synthetic quotes from \omterm{def:rc:sabr-in-rates-and-the-volatility-cube:normal}{normal SABR}; the chapter's
tutorial.}\label{fig:rc:sabr-in-rates-and-the-volatility-cube:neg}
\end{omfigure}

\section{Building the cube smile by smile}

\begin{definition}[Swaption matrix]\label{def:rc:sabr-in-rates-and-the-volatility-cube:matrix}
The \emph{swaption matrix}\index{swaption matrix} of a currency is the table of
at-the-money swaption volatilities by option expiry and underlying swap tenor:
the at-the-money face of the volatility cube (One Quant Book 2, chapter 13).
\end{definition}

\begin{method}[Calibrating the cube]\label{met:rc:sabr-in-rates-and-the-volatility-cube:cube}
For each quoted (expiry, tenor): compute the forward swap rate and its annuity
on the multi-curve (chapter~2); convert the quotes (normal volatilities, or
prices) into the model's own volatility (shifted Black volatilities for \omterm{def:rc:sabr-in-rates-and-the-volatility-cube:shifted}{shifted
SABR}); fix $\beta$ and the shift for the whole currency; fit $(\alpha,\rho,\nu)$ by
least squares, with $\alpha>0$ and $|\rho|<1$ imposed by reparametrisation;
record the fit error. Then check the parameters for smoothness across the grid:
a section whose $\rho$ or $\nu$ jumps is usually a bad quote.
\end{method}

\begin{omfigure}
\begin{tikzpicture}[x={(1cm,0cm)},y={(0.5cm,0.4cm)},z={(0cm,1cm)},font=\small]
\draw[omDef,thick] (0,0,0) -- (5,0,0) -- (5,4,0) -- (0,4,0) -- cycle;
\draw[omDef,thick] (0,0,3) -- (5,0,3) -- (5,4,3) -- (0,4,3) -- cycle;
\foreach \x/\y in {0/0,5/0,5/4,0/4} {\draw[omDef,thick] (\x,\y,0) -- (\x,\y,3);}
\fill[omThm!18] (0,0,1.5) -- (5,0,1.5) -- (5,4,1.5) -- (0,4,1.5) -- cycle;
\draw[omThm,thick] (0,0,1.5) -- (5,0,1.5) -- (5,4,1.5) -- (0,4,1.5) -- cycle;
\draw[omProp,very thick] (2.5,2,0) -- (2.5,2,3);
\node[anchor=north] at (2.5,0,0) {option expiry};
\node[anchor=west] at (5.1,2.2,0) {swap tenor};
\node[anchor=east] at (0,0,1.5) {strike};
\node[text=omThm,anchor=west] (atm) at (5.9,4,1.1) {at-the-money face: the swaption matrix};
\node[text=omProp,anchor=west] (sec) at (5.9,4,3.0) {one smile section: one SABR fit};
\draw[omThm,thin] (atm.west) -- (4.2,3.2,1.5);
\draw[omProp,thin] (sec.west) -- (2.5,2,2.6);
\end{tikzpicture}

\omcaption{The volatility cube: a smile for each (expiry, tenor). The
at-the-money face is the \omterm{def:rc:sabr-in-rates-and-the-volatility-cube:matrix}{swaption matrix}; each vertical line is a section
calibrated with its own SABR parameters, $\beta$ and the shift being fixed for the
whole currency.}\label{fig:rc:sabr-in-rates-and-the-volatility-cube:cube}
\end{omfigure}

\begin{example}[A synthetic euro cube]\label{ex:rc:sabr-in-rates-and-the-volatility-cube:cube}
Sixteen sections (expiries one, two, five and ten years; tenors two, five, ten
and thirty years), seven strikes each from $-200$ to $+200$ basis points around
the forward on chapter~2's curves, generated by \omterm{def:rc:sabr-in-rates-and-the-volatility-cube:normal}{normal SABR} with parameters
that vary smoothly with expiry and tenor plus seeded noise of 0.3 basis point,
are calibrated with \omterm{def:rc:sabr-in-rates-and-the-volatility-cube:shifted}{shifted SABR} ($\beta=0.5$, shift 3\%). Every section fits to
between 0.21 and 0.78 basis point. The at-the-money face
(\cref{fig:rc:sabr-in-rates-and-the-volatility-cube:matrix}) rises with expiry and
falls with tenor: 69.9 basis points for one year into two, 74.0 for ten into
thirty.
\end{example}

\begin{omfigure}
\begin{tikzpicture}
\begin{axis}[width=11cm,height=5.4cm,
  xlabel={option expiry (years)},ylabel={ATM normal volatility (bp)},
  tick label style={font=\small},label style={font=\small},
  xmin=0,xmax=10.5,ymin=50,ymax=95,grid=major,grid style={gray!25},
  legend columns=4,legend style={font=\footnotesize,at={(0.5,-0.25)},anchor=north,draw=none,fill=none,/tikz/every even column/.append style={column sep=6pt}},
  table/col sep=comma]
\addplot[omDef,thick,mark=*] table[x=e,y=n2]{figdata/rates-credit-risk/05-sabr-in-rates-and-the-volatility-cube/matrix.csv};
\addplot[omThm,thick,mark=square*] table[x=e,y=n5]{figdata/rates-credit-risk/05-sabr-in-rates-and-the-volatility-cube/matrix.csv};
\addplot[omProp,thick,mark=triangle*] table[x=e,y=n10]{figdata/rates-credit-risk/05-sabr-in-rates-and-the-volatility-cube/matrix.csv};
\addplot[omExo,thick,mark=diamond*] table[x=e,y=n30]{figdata/rates-credit-risk/05-sabr-in-rates-and-the-volatility-cube/matrix.csv};
\legend{2-year tenor,5-year,10-year,30-year}
\end{axis}
\end{tikzpicture}

\omcaption{The \omterm{def:rc:sabr-in-rates-and-the-volatility-cube:matrix}{swaption matrix} of the synthetic cube: at-the-money normal
volatility of the calibrated sections by expiry, one line per tenor. Data: the
chapter's tutorial.}\label{fig:rc:sabr-in-rates-and-the-volatility-cube:matrix}
\end{omfigure}

\section{Interpolating in expiry, tenor and strike}

The cube is quoted on a sparse grid and needed everywhere: a seven-year swap
starting in three years, a swaption struck anywhere. The strike dimension is
handled by the SABR formula itself; expiry and tenor are handled by
interpolating the parameters, not the volatilities, so that every interpolated
smile is itself a SABR smile.

\begin{example}[A section nobody quotes]\label{ex:rc:sabr-in-rates-and-the-volatility-cube:missing}
Interpolating $(\alpha,\rho,\nu)$ bilinearly in expiry and tenor between the
two-, five-year expiries and five-, ten-year tenors gives the three-year into
seven-year smile (forward 2.905\%) to within 0.93 basis point of the smile the
generating parameters imply, over the seven strikes.
\end{example}

\begin{remark}[What interpolation should respect]\label{rem:rc:sabr-in-rates-and-the-volatility-cube:interp}
Linear interpolation of $\alpha$ in expiry is a crude rule: total variance, not
volatility, should increase with expiry, and $\alpha$ is only the leading term.
Desks interpolate in variance along expiry, keep $\rho$ and $\nu$ smooth, and
check the interpolated sections for calendar and butterfly arbitrage (One Quant
Book 5, chapter 7). Forwards come from the curve, never from interpolation.
\end{remark}

\section{Arbitrage in the wings}

Hagan's formula is an expansion, accurate for moderate $\nu^2T$ and strikes not
too far from the forward. Its implied density, the second strike derivative of
the call price (Breeden--Litzenberger, One Quant Book 5, chapter 7), can become
negative in the low-strike wing of long expiries: a butterfly with negative
price.

\omcode{code/firm/sabrcube/firm_sabrcube.py}{136}{140}{The implied density of a section, by second differences of prices.}\label{lst:rc:sabr-in-rates-and-the-volatility-cube:density}

\begin{example}[Ten years into ten]\label{ex:rc:sabr-in-rates-and-the-volatility-cube:density}
A ten-year option on a ten-year swap at 2.50\%, quoted from $-200$ to $+200$
basis points around the forward (86 to 96 basis points of normal volatility),
is fitted by \omterm{def:rc:sabr-in-rates-and-the-volatility-cube:shifted}{shifted SABR} with a shift of 1\% (error 1.28 basis points) and of
3\% (0.35). With the 1\% shift the density is negative from $-0.95\%$ to
$+0.35\%$, reaching $-158$; with 3\% it stays positive down to $-2.15\%$
(\cref{fig:rc:sabr-in-rates-and-the-volatility-cube:density}). A book of
low-strike receivers priced with the 1\% fit carries an arbitrage against the
market's butterflies.
\end{example}

\begin{omfigure}
\begin{tikzpicture}
\begin{axis}[width=12cm,height=5.6cm,
  xlabel={rate at expiry (\%)},ylabel={density},
  tick label style={font=\small},label style={font=\small},
  xmin=-1,xmax=7,ymin=-165,ymax=30,grid=major,grid style={gray!25},unbounded coords=jump,
  legend columns=3,legend style={font=\footnotesize,at={(0.5,-0.25)},anchor=north,draw=none,fill=none,/tikz/every even column/.append style={column sep=8pt}},
  table/col sep=comma]
\addplot[omThm,thick] table[x=k,y=s1]{figdata/rates-credit-risk/05-sabr-in-rates-and-the-volatility-cube/density.csv};
\addplot[omDef,thick,dashed] table[x=k,y=s3]{figdata/rates-credit-risk/05-sabr-in-rates-and-the-volatility-cube/density.csv};
\addplot[omProp,thick,dotted] table[x=k,y=nrm]{figdata/rates-credit-risk/05-sabr-in-rates-and-the-volatility-cube/density.csv};
\legend{shifted SABR 1\%,shifted SABR 3\%,normal SABR (generating model)}
\draw[gray] (axis cs:-1,0) -- (axis cs:7,0);
\end{axis}
\end{tikzpicture}

\omcaption{Implied densities of the ten-year into ten-year swap rate from
three fits to the same quotes. The 1\% shift's density goes deeply negative
just above its floor; the 3\% shift's stays positive over the range shown.
Data: synthetic quotes from \omterm{def:rc:sabr-in-rates-and-the-volatility-cube:normal}{normal SABR}; the chapter's
tutorial.}\label{fig:rc:sabr-in-rates-and-the-volatility-cube:density}
\end{omfigure}

\begin{remark}[Fixes for the wings]\label{rem:rc:sabr-in-rates-and-the-volatility-cube:free}
Three are used: a larger shift (moving the problem below any relevant strike);
solving the SABR dynamics exactly or numerically instead of expanding (the
density is then positive by construction); and the free-boundary SABR of
Antonov, Konikov and Spector (2015), which replaces $F^\beta$ by $|F|^\beta$ and
lets the rate cross zero without a shift, with an exact price for zero
correlation. Whatever the model, a density check over the strikes a book
contains is a standard validation test (\cref{ch:rc:model-risk-and-validation}).
\end{remark}

\section{Tutorial: a cube from quotes}\label{tut:rc:sabr-in-rates-and-the-volatility-cube:tutorial}

\begin{tutorial}
\textbf{Goal.} Fit a negative-rate smile with three models, calibrate and
interpolate a cube, and check its density. \textbf{End state:}
\cref{fig:rc:sabr-in-rates-and-the-volatility-cube:neg,fig:rc:sabr-in-rates-and-the-volatility-cube:matrix,fig:rc:sabr-in-rates-and-the-volatility-cube:density}.
\begin{tutsteps}
\item \textbf{One smile}: \texttt{neg\_fits()} calibrates \omterm{def:rc:sabr-in-rates-and-the-volatility-cube:shifted}{shifted SABR} at 1\% and
3\% and \omterm{def:rc:sabr-in-rates-and-the-volatility-cube:normal}{normal SABR}; \texttt{deep\_receiver()} prices the $-0.90\%$ receiver.
\item \textbf{The cube}: \texttt{build\_cube()} returns the calibrated sections and
their errors; \texttt{matrix\_table()} the at-the-money face.
\item \textbf{Interpolate}: \texttt{missing\_section(3, 7)}.
\item \textbf{Density}: \texttt{density\_table()}; \texttt{fig\_rc\_sabrcube.py}
writes the charts.
\end{tutsteps}
\textbf{What to change next.} Fit with $\beta=0$ and $\beta=1$ at a 3\% shift and
compare $\rho$; interpolate total variance instead of $\alpha$ along expiry and
measure the error on the missing section.
\end{tutorial}

\section{Build: the swaption cube}\label{bld:rc:sabr-in-rates-and-the-volatility-cube:sabrcube}

\begin{build}
\bfield{Purpose} The volatility surface of every rates option of the book:
swaptions of any expiry, tenor and strike, CMS replication (chapter~6),
Bermudan calibration targets (chapters 7 to 9).
\bfield{Interface} \texttt{Sabr(alpha, beta, rho, nu, shift)};
\texttt{hagan\_lognormal}, \texttt{hagan\_normal\_beta0}, \texttt{price},
\texttt{normal\_vol}; \texttt{calibrate\_section(f, t, strikes, normal\_vols,
beta, shift, model)}; \texttt{density}; \texttt{Cube(expiries, tenors, sections,
forwards)} with \texttt{params}, \texttt{forward}, \texttt{normal\_vol}.
\bfield{Rules} Quotes in normal volatility; $\beta$ and the shift fixed per
currency; parameters interpolated, forwards taken from the curve; numpy only.
\bfield{Acceptance tests} \texttt{code/firm/sabrcube/tests/}: the at-the-money
limit of the expansion; \omterm{def:rc:sabr-in-rates-and-the-volatility-cube:normal}{normal SABR} without vol-of-vol is Bachelier; the
calibrator recovers planted parameters; payer minus receiver is the forward;
a benign density is positive and integrates to one; the cube interpolates
parameters bilinearly.
\bfield{Stretch} Interpolation in total variance; a PDE or exact-integration
SABR for positive densities; calibration to prices with bid--offer weights.
\end{build}

\begin{omsources}
\item P.\ Hagan, D.\ Kumar, A.\ Lesniewski and D.\ Woodward, ``Managing smile
risk'', \emph{Wilmott Magazine}, 2002.
\item A.\ Antonov, M.\ Konikov and M.\ Spector, ``The free boundary SABR:
natural extension to negative rates'', \emph{Risk}, 2015.
\item ECB, \emph{Key ECB interest rates}.
\item L.\ Andersen and V.\ Piterbarg, \emph{Interest Rate Modeling}, volume 2,
Atlantic Financial Press, 2010, chapter 16.
\end{omsources}

\section{Exercises}

\begin{exercise}[$\star$]\label{exo:rc:sabr-in-rates-and-the-volatility-cube:1}
A forward is $-0.30\%$ and a strike $-0.50\%$. Can standard lognormal SABR price
the option? What is the smallest shift that makes it possible?
\end{exercise}

\begin{exercise}[$\star$]\label{exo:rc:sabr-in-rates-and-the-volatility-cube:2}
In \omterm{def:rc:sabr-in-rates-and-the-volatility-cube:normal}{normal SABR}, the forward rises by 20 basis points overnight and the
parameters are unchanged. What happens to the at-the-money normal volatility,
and to the volatility of the strike that was at the money yesterday?
\end{exercise}

\begin{exercise}[$\star$]\label{exo:rc:sabr-in-rates-and-the-volatility-cube:3}
Why do rates desks fix $\beta$ instead of fitting it?
\end{exercise}

\begin{exercise}[$\star\star$]\label{exo:rc:sabr-in-rates-and-the-volatility-cube:4}
With \omterm{def:rc:sabr-in-rates-and-the-volatility-cube:normal}{normal SABR} at $\alpha=25$ basis points, $\rho=0.1$, $\nu=0.6$ and $T=1$,
compute the at-the-money normal volatility.
\end{exercise}

\begin{exercise}[$\star\star$]\label{exo:rc:sabr-in-rates-and-the-volatility-cube:5}
In \cref{ex:rc:sabr-in-rates-and-the-volatility-cube:neg}, why is the 1\% shift
fit worse at the quoted strikes as well as different in the wing?
\end{exercise}

\begin{exercise}[$\star\star$]\label{exo:rc:sabr-in-rates-and-the-volatility-cube:6}
Read \cref{fig:rc:sabr-in-rates-and-the-volatility-cube:density}. Describe the
butterfly that would exploit the 1\% fit near 0\%, and why it is not free money
in practice.
\end{exercise}

\begin{exercise}[$\star\star\star$]\label{exo:rc:sabr-in-rates-and-the-volatility-cube:7}
\emph{Coding.} Among shifts of 0.85\%, 0.90\%, 0.95\% and 1\%, find the smallest
for which the calibrated one-year-into-two smile has a positive density from
$-0.80\%$ to $+1\%$.
\end{exercise}

\begin{exercise}[$\star\star\star$]\label{exo:rc:sabr-in-rates-and-the-volatility-cube:8}
\emph{Find the flaw.} ``To fill the three-into-seven section we interpolate the
quoted normal volatilities strike by strike between the neighbouring sections.''
What goes wrong?
\end{exercise}

\section{Problem: The Minus-Ninety Strike}

\begin{problem}[{Weekend problem --- a client asks for a deep receiver}]\label{pb:rc:sabr-in-rates-and-the-volatility-cube:1}
In a negative-rate year a pension fund asks a desk for a receiver swaption:
one year into two years, on EUR~500 million, struck at $-0.90\%$, to protect
against rates falling much further. The forward is $-0.30\%$, the annuity
2.0, and the market quotes the six strikes of
\cref{ex:rc:sabr-in-rates-and-the-volatility-cube:neg}.

\medskip\noindent\textbf{Part I --- The quotes.}
\begin{enumerate}
\item How far is the strike from the forward, and from the lowest quote?
\item Why can standard lognormal SABR not price it?
\item Which fits are possible, and what does each assume about the lowest rate?
\item Give the fit errors of the three models.
\item Which is the best fit at the quoted strikes?
\end{enumerate}

\medskip\noindent\textbf{Part II --- The prices.}
\begin{enumerate}[resume]
\item Give the premium in euros under each model.
\item By how much do the 1\% and 3\% shifts disagree?
\item Why is the 1\% shift's price the lowest?
\item Which price would the desk rather sell at, and which would the client
rather buy at?
\item What is the smallest shift whose density is positive from $-0.80\%$ up?
\end{enumerate}

\medskip\noindent\textbf{Part III --- Risk.}
\begin{enumerate}[resume]
\item The desk sells at the 3\% price and hedges with the quoted $-0.80\%$
receiver. What risk remains?
\item How does the answer change if the ECB cuts to $-0.75\%$?
\item Why is the shift itself a model parameter that needs a reserve?
\item What does \omterm{def:rc:sabr-in-rates-and-the-volatility-cube:normal}{normal SABR} assume about how low rates can go?
\item How would you test the chosen model before quoting?
\end{enumerate}

\medskip\noindent\textbf{Part IV --- Judgement.}
\begin{enumerate}[resume]
\item Why did the euro market settle on shifts of a few per cent rather than
the smallest that works?
\item Is a single shift for the whole cube reasonable?
\item What would you tell the client about the price range?
\item State the \emph{named result}: the premium under shifts of 1\% and 3\%, and
the smallest admissible shift.
\item In one sentence: what decides the price of a strike nobody quotes?
\end{enumerate}
\end{problem}

\section{Interview questions}

\begin{interviewq}[$\star$ \iqroles{trader, researcher}]\label{iq:rc:sabr-in-rates-and-the-volatility-cube:1}
How did the rates market price options when rates went negative?
\end{interviewq}

\begin{interviewq}[$\star\star$ \iqroles{researcher}]\label{iq:rc:sabr-in-rates-and-the-volatility-cube:2}
What does $\beta$ control in SABR, and why is it usually fixed?
\end{interviewq}

\begin{interviewq}[$\star\star$ \iqroles{researcher, developer}]\label{iq:rc:sabr-in-rates-and-the-volatility-cube:3}
How would you build a swaption cube from broker quotes?
\end{interviewq}

\begin{interviewq}[$\star\star$ \iqroles{trader}]\label{iq:rc:sabr-in-rates-and-the-volatility-cube:4}
The forward moves up ten basis points. What happens to the smile under \omterm{def:rc:sabr-in-rates-and-the-volatility-cube:normal}{normal
SABR} and under lognormal SABR with $\beta=1$?
\end{interviewq}

\begin{interviewq}[$\star\star\star$ \iqroles{researcher, risk}]\label{iq:rc:sabr-in-rates-and-the-volatility-cube:5}
Hagan's formula gives a negative density. Where, why, and what do you do?
\end{interviewq}

\begin{interviewq}[$\star\star\star$ \iqroles{developer}]\label{iq:rc:sabr-in-rates-and-the-volatility-cube:6}
Your cube calibration fails on one section every morning. How do you
investigate?
\end{interviewq}
